\chapter{Introduction} \label{chap:intro}

In my project, I...

\section{Instruction Level Parallelism}

\textit{Instruction level parallelism} is the simultaneous execution of instructions in a computer processor. These 
instructions come from a single specific thread of execution in a computer program, as opposed to 
concurrency, where instructions from many threads are executed concurrently, and therefore is known as
\textit{thread level parallelism}.

Theoretically, a pair of instructions from a single thread can be executed simultaneously when there exist no data 
dependencies between them. A data dependency between two instructions exists when the operation of the instruction
appearing later in the program order consumes a value that is produced by the operation of the instruction appearing 
before it in the program order (this is known as a read after write, or RaW dependency).

To avoid race conditions, which lead to unexpected or inconsistent results, a computer processor has to take these 
RaW dependencies into consideration when scheduling instructions for execution. An instruction dependent on a previous
one must be executed strictly after the former has produced its result, to ensure correct propagation of values in the 
program.

\subsection{Finding and exploiting ILP in processors}

Previously, processors were in-order and single issue, having a maximum IPC of 1. Each clock cycle, the processor would 
execute one instruction, in program order, usually through a pipelined architecture.
Modern processors commonly exploit ILP, through superscalar processing and out of order execution.
Superscalar processors issue several independent instructions in a single clock cycle, commonly exhibiting an IPC 
greater than 1. This means that instructions are scheduled and executed out of order, which is a technique also used
to hide latencies from complex operations and memory loads.

In fact, the absolute limit of the ILP a processor core can exploit is entirely dependent on the structure of 
the computer program it executes. The dynamic execution of a program will have a 'critical path', that is, a series of 
executed instructions, each data-dependent on the previous, that must be executed sequentially. A processor cannot execute 
the program faster than it can execute this critical path.

However, the amount of ILP a processor is able to find and exploit is usually lower than the theoretical limit, owing to 
hardware limitations. These can include:

\begin{itemize}
  \item Branch/jump mispredictions, which typically trigger a discard of all work after the point of misprediction.
  \item Stuctural hazards, such as busy functional units or slow memory accesses.
  \item The number of instructions a superscalar processor can issue at once.
  \item The size of the window used when looking for independent instructions to schedule.
  \item False data dependencies, known as Write-after-Write or Write-after-Read, which are a consequence of the limited 
    number of architectural registers. They are mitigated in hardware using register renaming to a larger set of physical
    registers hidden to the programmer, but ultimately we may not resolve all these hazards.
\end{itemize}

\section{ILP Limit Studies}

%TODO: CITATIONS??

Increasing the amount of ILP a processor can find and exploit is a crucial part of high-performance processor design,
and so studies are conducted to justify the effort put into this problem.

These studies look to measure the theoretical limit of exploitable ILP in a program, by measuring the length of its 
critical path. The maximum ILP available in a program is then the number of instructions executed in the dynamic 
execution of the program, divided by the length of the critical path.

\subsection{Methodology of limit studies}

Idealised models of program execution are commonly used to find the absolute maximum upper bound of ILP. The simulators 
are fed a trace of an actual program execution on real hardware or a functional simulator. Since this trace already 
provides the program's path of execution, the simulator does not need to compute the results of instructions, and so it
purely acts as a dependency graph analyser and critical path calculator.

Instructions have \textit{static} dependencies, where each individual instruction will read from and write to a fixed 
set of registers and memory locations at its execution. These dependencies for the whole program can be represented by 
its full dynamic dependency graph (DDG), which is a directed acyclic graph with one node per dynamically executed 
instruction in a program. The edges of the graph represent dependencies between dynamically executed instructions,
and so the length of the critical path is the length of the longest path through the program's DDG.

% TODO: include diagram for DDG?

Limit studies will also vary the parameters of the analysis, using perfect models that ignore all hardware restrictions
and try to find the theoretical limit of ILP available in a program, or instead introducing realistic bounds, such as
control flow models (imperfect branch prediction), and instruction window scaling, where ILP is only computed within a
smaller window of instructions (instead of the whole program), to identify how 'local' a program's ILP is.

\section{Motivations}

% Why do we want to develop a new methodology for limit studies?

%TODO: SOURCE??

Traditionally, simulation models use an algorithm that iterates over the execution trace, and at each step calculates 
the length of dependency paths ending at each register, eventually picking out the longest path at the end of the trace.
Because the dynamic trace inherently guarantees that a producer instruction appears before its consumer,
the trace acts as a topological sort of the DDG. Therefore, starting at the first instruction in the trace,
the algorithm amounts to a single-pass linear sweep over this topologically sorted graph to calculate the longest path.

The weakness of this approach is that the traversal of the topological sort of the DDG is inherently sequential and 
this makes it unscalable for traces containing trillions of instructions. It cannot be heavily parallelised, nor can it 
exploit memoisation of repetitions in program stucture such as loops.

%TODO: CITATION

To break this sequential bottleneck, Chadwick proposes a matrix-based approach. The idea is to represent the dependencies 
between each two instructions in the trace as a matrix, and use matrix multiplication in the max-tropical semiring to 
compute the longest path in the DDG. This approach becomes a problem of reducing a large expression of matrix multiplications.

The immediately obvious benefit is that the matrix expression is a divisible computation, due to the associativity of 
matrix multiplication. The multiplications can be carried out in any order, and intermediate results fully encode the
dependency information of the instructions they range over. This enables a range of optimisations, including
thread-level parallelisation, result memoisation, and order of operations based on heuristics.

One additional benefit of the divisibility of the matrix expression is that it makes instruction window scaling for 
limit studies much easier. Consider three windows of instructions: A, B, and C (which is comprised of windows A and B
consecutively). To compute ILP information for each window, the traditional approach must be executed 3 times, once for 
each window. However, the matrix-based approach will only need to run twice, on A and B, and the dependency 
information matrix for C can simply be obtained by multiplying the result matrices from A and B.

In this project, I will explore several approaches for performing Chadwick's matrix-based approach efficiently, 
considering various types of parallelism, including at the data, thread, and accelerator level. I will then benchmark 
and evaluate these approaches, demonstrating a performance improvement over the 'classical' technique for computing 
dependency chains.
